% GENERATED FILE. Edit canonical lessons, then run npm run generate:books.
% canonical-source-hash: fnv1a64:4f468c40538cc15a
% canonical-lessons: ZH-C57-fangzi, ZH-C57-fangjian, ZH-C57-zhuozi, ZH-C57-yizi, ZH-C57-zher, ZH-R57-review-53-57, ZH-R57-review-final-57

\chapter{At Home}
\label{ch:zh-at-home}
\hlchaptermodality{57}

\begin{chapteropening}
\textbf{By the end of this chapter:} \emph{I can name my house, my room and the things in it, and say here.}

Complete the last lesson of chapter 57: I can name my house, my room and the things in it, and say here.
\end{chapteropening}

\section[fángzi]{fángzi — house}

\label{lesson:ZH-C57-fangzi}



\begin{quote}
Before the new one: say the Mandarin for to run, then the Mandarin for to laugh.
\end{quote}

\subsection*{You'll want to know: fángzi}
\emph{fángzi} — \textquotedblleft{}house\textquotedblright{}.

\begin{quote}
\emph{zhè shì wǒ de fángzi} — This is my house.
\end{quote}

Its characters come later; for now it is a word for the ear.

\begin{cousinweb}
\emph{fángzi} is built from fáng 'room' and the noun ending zi, both inherited from Middle Chinese, \emph{bjang} and \emph{tsiX} (Guangyun, in Baxter's transcription).
\end{cousinweb}

\begin{grammarlens}[title={What you've built}]
One more word about people.
\end{grammarlens}

\subsection*{Your turn}
\begin{itemize}
\raggedright
  \item \emph{Say it:} \emph{fángzi}
  \item \emph{Say it:} \emph{fángzi}, once more
  \item \emph{Say it:} \emph{zhè shì wǒ de fángzi} — This is my house.
  \item \emph{Recall:} say the Mandarin for to run, then the Mandarin for to laugh, then say \emph{fángzi} again
\end{itemize}

\begin{tcolorbox}[breakable,colback=teal!4,colframe=teal!35!black,title={Before you move on}]
What does fángzi mean? (\textquotedblleft{}House\textquotedblright{}.) Say it once more.
\end{tcolorbox}
\section[fángjiān]{fángjiān — room}

\label{lesson:ZH-C57-fangjian}



\begin{quote}
Before the new one: say the Mandarin for to laugh, then the Mandarin for house.
\end{quote}

\subsection*{You'll want to know: fángjiān}
\emph{fángjiān} — \textquotedblleft{}room\textquotedblright{}.

\begin{quote}
\emph{wǒ de fángjiān hěn xiǎo} — My room is small.
\end{quote}

Its characters come later; for now it is a word for the ear.

\begin{cousinweb}
\emph{fángjiān} is built from fáng 'room' and jiān 'space between', both inherited from Middle Chinese, \emph{bjang} and \emph{kean} (Guangyun, in Baxter's transcription); the dictionary lists the second under an older form of the character, glossed 'a gap'.
\end{cousinweb}

\begin{grammarlens}[title={What you've built}]
One more word about people.
\end{grammarlens}

\subsection*{Your turn}
\begin{itemize}
\raggedright
  \item \emph{Say it:} \emph{fángjiān}
  \item \emph{Say it:} \emph{fángjiān}, once more
  \item \emph{Say it:} \emph{wǒ de fángjiān hěn xiǎo} — My room is small.
  \item \emph{Recall:} say the Mandarin for to laugh, then the Mandarin for house, then say \emph{fángjiān} again
\end{itemize}

\begin{tcolorbox}[breakable,colback=teal!4,colframe=teal!35!black,title={Before you move on}]
What does fángjiān mean? (\textquotedblleft{}Room\textquotedblright{}.) Say it once more.
\end{tcolorbox}
\section[zhuōzi]{zhuōzi — table}

\label{lesson:ZH-C57-zhuozi}



\begin{quote}
Before the new one: say the Mandarin for house, then the Mandarin for room.
\end{quote}

\subsection*{You'll want to know: zhuōzi}
\emph{zhuōzi} — \textquotedblleft{}table\textquotedblright{}.

\begin{quote}
\emph{wǒ de zhuōzi hěn dà} — My desk is big.
\end{quote}

Its characters come later; for now it is a word for the ear.

\begin{cousinweb}
\emph{zhuōzi} is built from zhuō and the noun ending zi. The Guangyun lists the first character as \emph{traewk} (Guangyun, in Baxter's transcription), an old way of writing a word for 'high'; the table sense is not given there, so how it got it is not settled.
\end{cousinweb}

\begin{grammarlens}[title={What you've built}]
One more word about people.
\end{grammarlens}

\subsection*{Your turn}
\begin{itemize}
\raggedright
  \item \emph{Say it:} \emph{zhuōzi}
  \item \emph{Say it:} \emph{zhuōzi}, once more
  \item \emph{Say it:} \emph{wǒ de zhuōzi hěn dà} — My desk is big.
  \item \emph{Recall:} say the Mandarin for house, then the Mandarin for room, then say \emph{zhuōzi} again
\end{itemize}

\begin{tcolorbox}[breakable,colback=teal!4,colframe=teal!35!black,title={Before you move on}]
What does zhuōzi mean? (\textquotedblleft{}Table\textquotedblright{}.) Say it once more.
\end{tcolorbox}
\section[yǐzi]{yǐzi — chair}

\label{lesson:ZH-C57-yizi}



\begin{quote}
Before the new one: say the Mandarin for room, then the Mandarin for table.
\end{quote}

\subsection*{You'll want to know: yǐzi}
\emph{yǐzi} — \textquotedblleft{}chair\textquotedblright{}.

\begin{quote}
\emph{zhè shì háizi de yǐzi} — This is the child's chair.
\end{quote}

Its characters come later; for now it is a word for the ear.

\begin{cousinweb}
\emph{yǐzi} is built from yǐ and the noun ending zi. The Guangyun records the first character as \emph{'jeX} (Guangyun, in Baxter's transcription), and in a second reading as the name of a tree; the chair sense is not given there, so how it got it is not settled.
\end{cousinweb}

\begin{grammarlens}[title={What you've built}]
One more word about people.
\end{grammarlens}

\subsection*{Your turn}
\begin{itemize}
\raggedright
  \item \emph{Say it:} \emph{yǐzi}
  \item \emph{Say it:} \emph{yǐzi}, once more
  \item \emph{Say it:} \emph{zhè shì háizi de yǐzi} — This is the child's chair.
  \item \emph{Recall:} say the Mandarin for room, then the Mandarin for table, then say \emph{yǐzi} again
\end{itemize}

\begin{tcolorbox}[breakable,colback=teal!4,colframe=teal!35!black,title={Before you move on}]
What does yǐzi mean? (\textquotedblleft{}Chair\textquotedblright{}.) Say it once more.
\end{tcolorbox}
\section[\texorpdfstring{\zh{这儿} (zhèr)}{这儿 (zhèr)}]{\zh{这儿} (zhèr) — here}

\label{lesson:ZH-C57-zher}



\begin{quote}
Before the new one: say the Mandarin for table, then the Mandarin for chair.
\end{quote}

\subsection*{You'll want to know: \zh{这儿}}
\textbf{\zh{这儿}} — \emph{zhèr} — \textquotedblleft{}here\textquotedblright{}.

\begin{quote}
\textbf{\zh{你想来这儿吗}} \emph{nǐ xiǎng lái zhèr ma} — Would you like to come here?
\end{quote}

\begin{cousinweb}
\textbf{\zh{这儿}} (\emph{zhèr}) is built from \zh{这} 'this' and \zh{儿}, the ending that runs into the syllable before it as an r.
\end{cousinweb}

\begin{grammarlens}[title={What you've built}]
One more word about people.
\end{grammarlens}

\subsection*{Your turn}
\begin{itemize}
\raggedright
  \item \emph{Say it:} \emph{zhèr}
  \item \emph{Say it:} \emph{zhèr}, once more
  \item \emph{Say it:} \emph{nǐ xiǎng lái zhèr ma} — Would you like to come here?
  \item \emph{Recall:} say the Mandarin for table, then the Mandarin for chair, then say \emph{zhèr} again
\end{itemize}

\begin{tcolorbox}[breakable,colback=teal!4,colframe=teal!35!black,title={Before you move on}]
What does \zh{这儿} mean? (\textquotedblleft{}Here\textquotedblright{}.) Say it once more.
\end{tcolorbox}
\section[(dialogue)]{Looking back: from Speaking and Hearing to At Home}

\label{lesson:ZH-R57-review-53-57}



\begin{quote}
Say the Mandarin for to say and spoken words.
\end{quote}

\subsection*{Your turn}
\begin{itemize}
\raggedright
  \item \emph{Say it:} to say, spoken words, to listen, to understand, slow
  \item \emph{Say it:} to learn, to write, to read aloud, to ask, lesson
  \item \emph{Say it:} to want to, to love, film, to sing, coffee
  \item \emph{Say it:} to play, to play ball, sport, to run, to laugh
  \item \emph{Say it:} house, room, table, chair, here
  \item \emph{Read it:} \zh{说} \zh{话} \zh{想} — each one aloud, then a word that holds it
\end{itemize}

\begin{tcolorbox}[breakable,colback=teal!4,colframe=teal!35!black,title={Before you move on}]
To say? (\textbf{\zh{说}}, \emph{shuō}.) Spoken words? (\textbf{\zh{话}}, \emph{huà}.) To listen? (\emph{tīng}.) To understand? (\emph{dǒng}.) Slow? (\emph{màn}.) To learn? (\emph{xuéxí}.) To write? (\emph{xiě}.) To read aloud? (\emph{dú}.) To ask? (\emph{wèn}.) Lesson? (\emph{kè}.) To want to? (\textbf{\zh{想}}, \emph{xiǎng}.) To love? (\emph{ài}.) Film? (\emph{diànyǐng}.) To sing? (\emph{chànggē}.) Coffee? (\emph{kāfēi}.) To play? (\emph{wán}.) To play ball? (\emph{dǎqiú}.) Sport? (\emph{yùndòng}.) To run? (\emph{pǎobù}.) To laugh? (\emph{xiào}.) House? (\emph{fángzi}.) Room? (\emph{fángjiān}.) Table? (\emph{zhuōzi}.) Chair? (\emph{yǐzi}.) Here? (\textbf{\zh{这儿}}, \emph{zhèr}.)
\end{tcolorbox}
\section[(dialogue)]{At Home — once more}

\label{lesson:ZH-R57-review-final-57}



\begin{quote}
Say the Mandarin for house and room.
\end{quote}

\subsection*{Your turn}
\begin{itemize}
\raggedright
  \item \emph{Say it:} house, room, table, chair, here
\end{itemize}

\begin{tcolorbox}[breakable,colback=teal!4,colframe=teal!35!black,title={Before you move on}]
House? (\emph{fángzi}.) Room? (\emph{fángjiān}.) Table? (\emph{zhuōzi}.) Chair? (\emph{yǐzi}.) Here? (\textbf{\zh{这儿}}, \emph{zhèr}.)
\end{tcolorbox}
